\section{Introduction}

Modern analytical placers make global placement tractable by optimizing
half-perimeter wirelength (HPWL) plus a density penalty. These cheap, smooth
proxies enable gradient-based tools such as DREAMPlace~\cite{dreamplace}.
However, the same objective choice also limits
what the placer can optimize. ChiPBench shows that placement-stage HPWL
correlates poorly with post-route quality on modern industrial designs,
including routed wirelength, routing congestion, worst negative slack (WNS),
total negative slack (TNS), and design-rule checking (DRC)
count~\cite{chipbench}. The analytical placer optimizes one target while the
completed flow depends on another.

Routing-aware placement has addressed this mismatch through two main changes
to the objective. The first adds hand-engineered symbolic terms on top of a
fixed placement objective. Rectangular Uniform wire DensitY (RUDY) estimates
routing demand from placement geometry~\cite{rudy}, and later congestion and
density controls build on the same idea. Terms of this kind sharpen routing
awareness but fix the objective form and require expert redesign for each
technology or benchmark family. The second inserts learned routing or timing
surrogates into placement. LaMPlace improves post-route WNS by 43.0 percent
and TNS by 30.4 percent~\cite{lamplace}. GOALPlace reduces DRC violations by
up to tenfold~\cite{goalplace}. RoutePlacer lowers routing congestion while
preserving routed wirelength~\cite{routeplacer}. These methods add downstream
awareness but embed part of the objective in black-box models that engineers
cannot inspect or edit. Neither change makes the differentiable placement
objective itself searchable. This leaves an open third path: direct search over
readable, differentiable objective forms that still run inside analytical
placement.

CoEvoP\&R searches this third path with LLM program evolution. FunSearch,
Eureka, LLM-SR, and AlphaEvolve show that LLMs can revise executable programs
using automatic evaluator feedback~\cite{funsearch,eureka,llmsr,alphaevolve}.
EvoPlace and VeoPlace bring foundation-model search into placement by evolving
optimizer components or guiding macro-placement
actions~\cite{evoplace,veoplace}. These systems search initialization,
optimizer, or action choices around the analytical objective. CoEvoP\&R
searches the objective itself: an LLM proposes complete symbolic replacement objectives,
and DREAMPlace evaluates each safe candidate inside the gradient path that
produces a placement. This embedding makes optimizer
compatibility part of the search space. Static admission, diagnostic autograd
checks, and runtime guards reject incompatible or unstable objectives. Because
full OpenROAD routing is too expensive for every candidate, scheduled
promotions return routed evidence to later proposals~\cite{openroad}. A proxy
audit gives cheap timing signals selection influence only after they are
non-redundant with HPWL and agree with downstream timing.

CoEvoP\&R makes four contributions.
\begin{itemize}
  \item \textbf{Searchable objectives.} CoEvoP\&R searches complete
  differentiable replacement objectives and discovers readable formulas that
  improve placement and post-GRT quality over the native DREAMPlace objective.
  \item \textbf{Optimizer-embedded evolution.} CoEvoP\&R inserts each LLM
  objective into the DREAMPlace gradient path through static, diagnostic, and
  runtime safety checks.
  \item \textbf{Proxy audit.} CoEvoP\&R gives cheap timing proxies
  selection influence only after non-redundancy and routed-agreement checks.
  \item \textbf{Deployable artifact.} Each CoEvoP\&R search returns a compact
  symbolic objective that engineers can inspect, edit, and plug into DREAMPlace.
\end{itemize}

On the primary Nangate45 panel, CoEvoP\&R reduces DREAMPlace HPWL by 16.9\%
and density overflow by 36.7\% on average, while improving OpenTimer
placement-RC WNS by 0.70 ns and TNS by 912 ns.
